% Ngân hàng bài tập — chương 4
% Chương được xác định bởi tên của tệp này.

\begin{exo}[Tích phân một dạng vi phân dọc theo một đường]{integrale-forme-differentielle}
\keywords{dạng vi phân, vi phân toàn phần, tiêu chuẩn Schwarz, tích phân đường}

Xét hai dạng vi phân
\[
\omega_1 = y\,\mathrm{d}x + x\,\mathrm{d}y,
\qquad
\omega_2 = y\,\mathrm{d}x,
\]
và ba đường nối gốc $O(0,0)$ với điểm $M(1,1)$: đường $\gamma_1$ gồm đoạn
ngang $O\to(1,0)$ rồi đoạn đứng $(1,0)\to M$; đường $\gamma_2$ gồm đoạn đứng
$O\to(0,1)$ rồi đoạn ngang $(0,1)\to M$; và đường $\gamma_3$ là đoạn thẳng
nối trực tiếp $O$ với $M$.

\begin{enumerate}
\item Chứng minh $\omega_1$ là dạng vi phân toàn phần bằng cách tìm hàm $f$ sao cho $\omega_1 = \mathrm{d}f$.
\item Tính $\int_{\gamma_1}\omega_1$ và $\int_{\gamma_2}\omega_1$ bằng cách tham số hóa từng đoạn. Suy ra lại kết quả mà không cần tính toán.
\item Áp dụng tiêu chuẩn Schwarz cho $\omega_2$. Có thể kết luận gì?
\item Tính $\int_{\gamma_1}\omega_2$, $\int_{\gamma_2}\omega_2$ và $\int_{\gamma_3}\omega_2$. Bình luận.
\end{enumerate}

\begin{indication}
Trên đoạn ngang $\mathrm{d}y = 0$, trên đoạn đứng $\mathrm{d}x = 0$:
mỗi tích phân quy về một tích phân một biến thông thường.
\end{indication}

\begin{solution}
\textbf{Câu 1.} $f(x,y) = xy$ thỏa mãn vì $\partial f/\partial x = y$ và $\partial f/\partial y = x$.

\textbf{Câu 2.} Với đường cong được tham số hóa bởi $t\mapsto\bigl(x(t),y(t)\bigr)$,
\[
\int_\gamma\omega_1
= \int \left[y(t)x'(t)+x(t)y'(t)\right],\mathrm{d}t.
\]
Đường $\gamma_1$ gồm hai đoạn. Trên đoạn thứ nhất, $x(t)=t$, $y(t)=0$;
trên đoạn thứ hai, $x(t)=1$, $y(t)=t$; trong cả hai trường hợp $t\in[0,1]$. Do đó
\[
\begin{aligned}
I_1=\int_{\gamma_1}\omega_1
&= \int_0^1(0\times1+t\times0)\,\mathrm{d}t
+\int_0^1(t\times0+1\times1)\,\mathrm{d}t\\
&=0+\int_0^1\mathrm{d}t=1.
\end{aligned}
\]
Với $\gamma_2$, đoạn thứ nhất có $x(t)=0$, $y(t)=t$, đoạn thứ hai có
$x(t)=t$, $y(t)=1$, và $t\in[0,1]$. Vì vậy
\[
\begin{aligned}
I_2=\int_{\gamma_2}\omega_1
&= \int_0^1(t\times0+0\times1)\,\mathrm{d}t
+\int_0^1(1\times1+t\times0)\,\mathrm{d}t\\
&=0+\int_0^1\mathrm{d}t=1.
\end{aligned}
\]
Hai giá trị trùng nhau. Có thể dự đoán điều đó mà không cần tính: vì
$\omega_1=\mathrm{d}f$ với $f(x,y)=xy$, tích phân chỉ phụ thuộc vào hai đầu mút,
\[
\int_\gamma\omega_1=f(M)-f(O)=f(1,1)-f(0,0)=1.
\]

\textbf{Câu 3.} Viết $\omega_2=A(x,y)\,\mathrm{d}x+B(x,y)\,\mathrm{d}y$.
Ở đây $A(x,y)=y$, $B(x,y)=0$. Nếu dạng này là toàn phần, tiêu chuẩn Schwarz đòi hỏi
\[
\frac{\partial A}{\partial y}=\frac{\partial B}{\partial x}.
\]
Nhưng $\partial A/\partial y=1$ còn $\partial B/\partial x=0$. Hai đạo hàm chéo
khác nhau, nên $\omega_2$ không phải dạng vi phân toàn phần. Tích phân của nó có thể
phụ thuộc vào đường nối $O$ và $M$, như phép tính sau xác nhận.

\textbf{Câu 4.} Vì $\omega_2=y\,\mathrm{d}x$, chỉ biến thiên của $x$ tại $y\ne0$
mới có thể đóng góp vào tích phân. Với các tham số hóa trên,
\[
\begin{aligned}
J_1=\int_{\gamma_1}\omega_2
&=\int_0^1 0\times1\,\mathrm{d}t+\int_0^1 t\times0\,\mathrm{d}t=0.
\end{aligned}
\]
Trên đoạn đầu của $\gamma_2$, $x(t)=0$ nên $\mathrm{d}x=0$; trên đoạn sau,
$x(t)=t$, $y(t)=1$ và $\mathrm{d}x=\mathrm{d}t$. Do đó
\[
\begin{aligned}
J_2=\int_{\gamma_2}\omega_2
&=\int_0^1 t\times0\,\mathrm{d}t+\int_0^1 1\times1\,\mathrm{d}t=1.
\end{aligned}
\]
Cuối cùng, đoạn chéo $\gamma_3$ được tham số hóa bởi
\[
x(t)=t,\qquad y(t)=t,\qquad t\in[0,1],
\]
nên
\[
J_3=\int_{\gamma_3}\omega_2
=\int_0^1y(t)x'(t)\,\mathrm{d}t
=\int_0^1t\,\mathrm{d}t=\frac12.
\]
Ba đường có cùng đầu mút nhưng cho $J_1=0$, $J_2=1$, $J_3=1/2$ khác nhau.
Đó chính là tính phụ thuộc đường đi của một dạng vi phân không toàn phần.
\end{solution}
\end{exo}

\begin{exo}[Cùng biến thiên nội năng, ba phương thức truyền]{meme-delta-u-trois-transferts}
\keywords{nguyên lý thứ nhất, nội năng, nhiệt, công cơ học, công điện, phép đo nhiệt lượng, hàm trạng thái}

Một chất lỏng, nhiệt lượng kế của nó và một điện trở nhỏ nhúng trong chất lỏng tạo thành
một hệ kín, đứng yên ở cấp vĩ mô. Bình cứng và tổng nhiệt dung của hệ được coi là không đổi:
\[
C=2{,}00\ \mathrm{kJ\,K^{-1}}.
\]
Ta muốn đưa hệ từ trạng thái cân bằng $A$ ở $T_A=20{,}0\,^\circ\mathrm{C}$
đến trạng thái cân bằng $B$ ở $T_B=25{,}0\,^\circ\mathrm{C}$. Cho rằng
\[
U(B)-U(A)=C(T_B-T_A).
\]
Các biến thiên động năng và thế năng vĩ mô đều bằng không. Bỏ qua mọi thất thoát ra ngoài.

Ba quy trình sau được thực hiện độc lập, cùng bắt đầu từ trạng thái $A$:
\begin{itemize}
\item[Quy trình 1] Cho hệ tiếp xúc nhiệt với một nguồn nóng mà không cung cấp công cho hệ.
\item[Quy trình 2] Các thành được cách nhiệt. Một cánh khuấy khuấy chất lỏng nhờ vật khối lượng $m$ rơi qua độ cao $h=5{,}00$ m. Cuối cùng cánh khuấy và chất lỏng đứng yên; toàn bộ cơ năng mà vật mất đi được truyền cho hệ. Lấy $g=9{,}81\ \mathrm{m\,s^{-2}}$.
\item[Quy trình 3] Hệ nhận nhiệt $Q_3=4{,}00$ kJ, phần năng lượng còn thiếu được cấp cho điện trở bằng điện. Công suất điện nhận được không đổi, $\mathcal P=300$ W.
\end{itemize}

\begin{enumerate}
\item Tính biến thiên nội năng $\Delta U=U(B)-U(A)$ chung cho cả ba quy trình.
\item Với hai quy trình đầu, xác định nhiệt $Q$ và công $W$ mà hệ nhận được. Trong quy trình 2, tính khối lượng $m$ cần thiết.
\item Với quy trình 3, tính công điện nhận được và thời gian cấp điện cho điện trở.
\item Ba quá trình có cùng trạng thái đầu và cuối. Tại sao các giá trị $Q$ và $W$ vẫn có thể khác nhau? Ở trạng thái cuối $B$, hệ có chứa “nhiệt” hay “công” không?
\end{enumerate}

\begin{indication}
Áp dụng $\Delta U=Q+W$ theo quy ước dấu của nhà ngân hàng. Công hệ nhận qua cánh khuấy
bằng độ giảm thế năng $mgh$ của vật. Năng lượng điện đi qua biên theo dây dẫn được tính là công điện.
\end{indication}

\begin{solution}
\textbf{Câu 1.} Độ chênh nhiệt độ là $T_B-T_A=5{,}0$ K. Do đó
\[
\Delta U=C(T_B-T_A)=2{,}00\times5{,}0=10{,}0\ \mathrm{kJ}.
\]
Giá trị này chỉ phụ thuộc vào hai trạng thái cân bằng $A$ và $B$.

\textbf{Câu 2.} Trong quy trình 1, hệ không nhận công: $W_1=0$. Nguyên lý thứ nhất cho
\[
Q_1=\Delta U=10{,}0\ \mathrm{kJ}.
\]
Nhiệt dương vì năng lượng đi vào hệ. Trong quy trình 2, thành cách nhiệt nên $Q_2=0$.
Toàn bộ biến thiên nội năng đến từ công cơ học của cánh khuấy:
\[
W_2=\Delta U=10{,}0\ \mathrm{kJ}.
\]
Vì $W_2=mgh$,
\[
m=\frac{W_2}{gh}
=\frac{10{,}0\times10^3}{9{,}81\times5{,}00}
\simeq204\ \mathrm{kg}.
\]
Bậc độ lớn cao cho thấy một mức tăng nhiệt độ nhỏ đã tương ứng với cơ năng đáng kể.

\textbf{Câu 3.} Nguyên lý thứ nhất đòi hỏi
\[
W_3=\Delta U-Q_3=10{,}0-4{,}00=6{,}00\ \mathrm{kJ}.
\]
Công dương vì nguồn điện cung cấp năng lượng cho hệ. Với $W_3=\mathcal P\,\Delta t$,
\[
\Delta t=\frac{W_3}{\mathcal P}=\frac{6{,}00\times10^3}{300}=20{,}0\ \mathrm{s}.
\]

\textbf{Câu 4.} Ba đường lần lượt cho
\[
(Q_1,W_1)=(10{,}0,0)\ \mathrm{kJ},\qquad
(Q_2,W_2)=(0,10{,}0)\ \mathrm{kJ},
\]
và
\[
(Q_3,W_3)=(4{,}00,6{,}00)\ \mathrm{kJ}.
\]
Trong mọi trường hợp, $Q+W=10{,}0$ kJ và tạo cùng $\Delta U$. Nội năng là hàm trạng thái,
còn nhiệt và công mô tả các truyền năng lượng trong quá trình. Ở trạng thái cuối $B$, hệ có nội năng,
nhưng không chứa “nhiệt” hay “công”.
\end{solution}
\end{exo}

\begin{exo}[Nén dưới áp suất ngoài không đổi]{compression-pression-exterieure-constante}
\seoready{true}
\keywords{công của lực áp suất, áp suất ngoài không đổi, nén, giãn nở, giãn nở trong chân không}

Một chất khí bị nén từ thể tích $V_A$ đến thể tích $V_B<V_A$ dưới áp suất ngoài không đổi $P_0$.

\begin{enumerate}
\item Tính công mà chất khí nhận được.
\item Chất khí trở về từ $V_B$ đến $V_A$ chống lại cùng áp suất ngoài $P_0$. Tính công nhận được và so sánh với công nén.
\item Xét lại sự giãn nở từ $V_B$ đến $V_A$ khi khí giãn nở trong chân không. Giải thích ba dấu thu được.
\end{enumerate}

\begin{indication}
Dùng $W=-\int P_{\mathrm{ext}}\,\mathrm{d}V$. Phải tích phân áp suất \emph{ngoài},
kể cả khi quá trình không thuận nghịch giả tĩnh.
\end{indication}

\begin{solution}
\textbf{Câu 1.} Áp suất ngoài không đổi, nên
\[
W_{A\to B}=-P_0(V_B-V_A)=P_0(V_A-V_B)>0.
\]
Chất khí nhận công khi bị nén.

\textbf{Câu 2.} Với quá trình trở về,
\[
W_{B\to A}=-P_0(V_A-V_B)<0.
\]
Hai công đối nhau vì hai quá trình diễn ra chống lại cùng một áp suất ngoài không đổi.

\textbf{Câu 3.} Trong chân không, $P_{\mathrm{ext}}=0$, do đó
\[
W_{B\to A}^{\mathrm{vide}}=0.
\]
Vì vậy giãn nở không tự động kéo theo công âm: phải thực sự đẩy lùi một ngoại lực.
\end{solution}
\end{exo}
